% Methodology -- first draft, written 2026-10-05 after thesis/style_guide_methodology_aida.md.
% Scope: what is done on the Herlev-Osterbro (CGPS) cohort. ImageCAS-X appears only as the
% development set on which outcome-free choices are locked. Deliberately not tied to the code in
% this repo: tools are named by role, so the text survives a change of segmentation model or tracer.
%
% STANDING CONSTRAINT: no exact figure from CGPS in this file until data access. Every cohort
% number is a placeholder [n]; every open decision is a % TODO next to the sentence it affects.
% Feature and model choices follow thesis/week5/risk_model_plan.tex.

\chapter{Methodology}
\label{ch:methodology}

The methodology followed in this work is summarized in Figure~\ref{fig:method-overview} and
consists of six main stages: (1) selection of the study population and definition of the outcome;
(2) segmentation of the coronary lumen in each baseline scan; (3) extraction of the coronary
centerline tree and its orientation from the ostium to the distal vessels; (4) labeling of each
centerline segment with the artery it belongs to; (5) computation of geometric features at the
bifurcations and along the arteries; and (6) statistical modeling of new plaque from these
features.
All choices in stages 2 to 5 that do not depend on the outcome were fixed on a public development
dataset before any follow-up data from Herlev--{\O}sterbro were opened.

\begin{figure}[htbp]
  \centering
  % TODO: overview figure, one box per stage, ImageCAS-X shown as a side lane feeding stages 2 to 5.
  \fbox{\parbox{0.9\linewidth}{\centering\vspace{2cm}Placeholder: methodology overview\vspace{2cm}}}
  \caption{Overview of the methodology, from the baseline scan to the risk model. The public
    development dataset is used to train the segmentation and labeling models and to fix every
    outcome-free parameter before the Herlev--{\O}sterbro outcomes are used.}
  \label{fig:method-overview}
\end{figure}


\section{Study Population and Outcome}
\label{sec:method-population}

The analysis uses the coronary computed tomography angiography (CCTA) sub-study of the
Copenhagen General Population Study, referred to in this thesis as Herlev--{\O}sterbro
\cite{fuchs2023cgps}.
Participants were scanned as part of a research protocol rather than on clinical indication, so
their coronary trees sample a general population rather than a referred one.
A subset of participants underwent a second CCTA scan, and this subset forms the study population.
% TODO: rescan protocol, selection of who was rescanned, and the scan interval once data access is granted.

Participants were included if their whole coronary tree was free of plaque at the baseline scan.
The restriction applies to the whole tree rather than to the artery being modeled, since an
artery enlarges outward as plaque accumulates \cite{glagov1987remodeling} and a tree with plaque
in one artery may already carry remodeled geometry in another.
Scans with insufficient image quality or incomplete coverage of the coronary tree were excluded.
From the [n] participants with two scans, [n] were retained after these criteria, and every
exclusion was recorded with its reason.
% TODO: exclusion counts per reason, as a flow diagram or a table.

The outcome is new plaque at the follow-up scan, defined separately for four artery domains: the
left main (LM), the left anterior descending (LAD) artery, the left circumflex (LCx) artery and
the right coronary artery (RCA).
A domain was counted as an event when plaque was present at follow-up in any of its segments.
% TODO: who graded plaque at follow-up, on which segment scheme (SCCT?), and whether calcified and
% non-calcified plaque are pooled.
Side branches were assigned to the domain of the artery they arise from.
Defining the outcome per domain rather than per participant keeps a harmful geometry in one artery
from being averaged with a harmless one in another, as wall shear stress acts locally
\cite{chatzizisis2007ess}.

Age, sex and coronary dominance were recorded for every participant, together with the classical
risk factors available in the study.
% TODO: list the risk factors and the calcium score definition available in CGPS.


\section{Lumen Segmentation}
\label{sec:method-segmentation}

The coronary lumen was segmented automatically in each baseline scan, since manual segmentation of
a population cohort is not feasible.
A deep-learning segmentation model was trained on ImageCAS-X \cite{bransby2026imagecasx}, a
curated release of the public ImageCAS dataset \cite{zeng2023imagecas} with corrected coronary
segmentations.
The model was chosen from the methods benchmarked on that dataset, and its published training
settings were followed.
% TODO: name the model and cite it. Candidates: CAS-Net \cite{dong2023casnet}, nnU-Net \cite{isensee2021nnunet}.
Every deviation from the published settings is listed in Table~\ref{tab:method-deviations} with
its reason.

ImageCAS-X was split into training, validation and test sets as provided with the dataset, and the
test set was used only to report segmentation performance.
Performance was measured with the Dice similarity coefficient between the predicted segmentation
$P$ and the reference $G$,
\begin{equation}
  \mathrm{Dice}(P, G) = \frac{2\,|P \cap G|}{|P| + |G|},
  \label{eq:method-dice}
\end{equation}
where $|\cdot|$ denotes the number of voxels in a region.
Dice measures overlap and is dominated by the large proximal vessels, so connectivity was also
measured with the centerline Dice (clDice) \cite{shit2021cldice}, which scores how much of each
centerline lies inside the other segmentation.
A break in a vessel changes Dice little but changes every feature measured downstream of it, which
makes clDice the more relevant of the two for this work.

The trained model was then applied to the Herlev--{\O}sterbro baseline scans without further
training, since no reference segmentations exist for this cohort.
Each scan was resampled to the voxel spacing used in training, and the prediction was resampled
back to the original geometry.
Connected components smaller than a fixed volume were removed, as they are not attached to the
coronary tree.
A random sample of [n] predictions was inspected visually in three orthogonal views and in 3D,
and the failure modes found were counted.
% TODO: automatic quality flags (number of components, total tree length, missing major artery)
% and the count of scans excluded on them.


\section{Centerline Extraction and Tree Construction}
\label{sec:method-tree}

The centerline of each segmented tree was extracted and represented as a set of points in physical
coordinates joined by polylines.
% TODO: state which extractor is used on CGPS: the group's in-house tracer or a skeleton of the
% segmentation. If the in-house tracer, cite it and say it was used without modification.
A centerline alone describes the geometry of the vessels but not the direction of flow through
them, so each tree was converted into a rooted graph oriented from the ostium to the distal
vessels.

Points where three or more polylines meet were taken as bifurcations and points where a single
polyline ends were taken as termini.
Polylines were joined through points shared by exactly two of them, so that each segment of the
graph runs from one junction to the next.
The ostium of each side was placed at the centerline point closest to the aortic surface, with
the aorta segmented by TotalSegmentator \cite{wasserthal2023totalsegmentator}.
Starting from the ostium, a breadth-first traversal oriented every segment away from the root and
recorded its parent segment, its child segments and its generation, defined as the number of
bifurcations between the segment and the ostium.

Three departures from a single rooted tree were handled explicitly rather than discarded.
A left tree without an LM, in which the LAD and LCx arise from the aorta through separate ostia
\cite{topaz1991absentlm}, was represented as two rooted trees.
A fragment disconnected from the tree was rooted at its end nearest to the rest of the tree, since
that end is the most plausible proximal point.
A loop in the graph is not anatomical in a coronary tree, so the segment closing it was removed
and the case was flagged.
The number of trees affected by each departure was recorded.


\section{Artery Labeling}
\label{sec:method-labeling}

The features in Section~\ref{sec:method-features} are defined per named artery, but the extracted
trees carry no artery names.
Each segment was therefore labeled with one of the coronary artery classes of ImageCAS-X by a
graph neural network operating on the oriented tree \cite{hampe2024gnn}.
The network was trained on the labeled centerline trees delivered with ImageCAS-X, with segment
position, direction, length and radius as input, expressed in a frame defined by the heart
chambers so that the input does not depend on how the patient lay in the scanner.
% TODO: final input features, the heart frame definition, and the label set once fixed.

Accuracy was measured on the ImageCAS-X test trees, both on the delivered centerlines and on
centerlines extracted from predicted segmentations, since only the second reflects the trees the
network sees in Herlev--{\O}sterbro.
For the participants with disease, Herlev--{\O}sterbro also provides artery labels produced by
commercial software (Medis).
These labels are generated automatically and were treated as a reference of uncertain quality, so
agreement with them was reported together with a visual check of [n] disagreements.
They exist only for participants with disease, whereas the study population is plaque free at
baseline, so labeling accuracy on the study population was estimated from a visual check of [n]
randomly drawn baseline trees.
% TODO: mapping from Medis segment names onto the ImageCAS-X label set.

Two labeling outcomes need a rule, as they occur in normal anatomy.
Trees with a ramus intermedius were flagged, and every analysis was repeated without them.
A missing first diagonal (D1) or first obtuse marginal (OM1) was coded as absent rather than
treated as missing data, since the absence is itself a property of the tree.


\section{Geometric Features}
\label{sec:method-features}

All features were computed from the oriented centerline and the vessel radius, without flow
simulation, since simulating flow through every tree of a population cohort is not feasible and
depends on assumed boundary conditions.
Each feature was chosen for a published link between the geometry it measures and wall shear
stress.
The tortuosity index, torsion and root-mean-square curvature were considered and not used; the
reasons are given in Section~\ref{sec:method-dropped}.

\subsection{Vessel Radius}

The radius at each centerline point was taken as the distance from that point to the lumen
boundary, computed with a Euclidean distance transform of the segmentation.
Near a bifurcation the lumens of the parent and daughter vessels merge, and a radius or direction
measured there describes the junction rather than either vessel.
Radii and directions at a bifurcation were therefore measured over a 3\,mm window of each vessel
that starts one junction radius beyond the branch point, following the inscribed-sphere convention
of VMTK \cite{antiga2008vmtk}.
The radius of a vessel at a bifurcation is the median over that window, which is less sensitive to
single noisy points than the mean.

\subsection{Bifurcation Features}

Three bifurcations were measured in each left tree: the LM into the LAD and LCx, the LAD into the
D1, and the LCx into the OM1.
At each, the direction of each daughter vessel was the principal axis of its window points,
oriented away from the junction, and the bifurcation angle is
\begin{equation}
  \theta = \arccos\left(\mathbf{d}_1 \cdot \mathbf{d}_2\right),
  \label{eq:method-angle}
\end{equation}
where $\mathbf{d}_1$ and $\mathbf{d}_2$ are the unit directions of the two daughters.

The balance of a bifurcation was described by the daughter ratio
\begin{equation}
  \rho = \frac{\min(r_1, r_2)}{\max(r_1, r_2)},
  \label{eq:method-rho}
\end{equation}
where $r_1$ and $r_2$ are the daughter radii.
At the LM, the outcome models for the LAD and LCx domains instead use the signed log ratio
$\lambda = \ln(r_\mathrm{LAD}/r_\mathrm{LCx})$, which keeps which daughter is larger.
$\rho$ and $\lambda$ carry the same information at the LM, so they never enter the same model.
Ratios of radii were preferred over radii, as they cancel any scale error common to both vessels
\cite{finet2008fractal}.

The relation between the parent radius $r_0$ and the daughter radii was compared with the scaling
law of Huo and Kassab \cite{huo2012scaling} through the deviation
\begin{equation}
  \Delta_{HK} = \frac{1}{r_0}\left(\frac{r_1^{7/3} + r_2^{7/3}}{2}\right)^{3/7} - 2^{-3/7},
  \label{eq:method-hk}
\end{equation}
which is zero for a bifurcation that follows the law.
No coronary study has related $\Delta_{HK}$ to plaque, so its direction of effect was treated as
unknown and tested two-sided.

\subsection{Curvature}

Curvature was measured along each named artery from its origin.
The centerline was smoothed before differentiation, as curvature depends on the second derivative
of the point positions and amplifies their noise.
% TODO: smoothing method and its parameter, fixed on ImageCAS-X.
For a centerline $\mathbf{c}(s)$ parametrized by arc length $s$, the curvature is
$\kappa(s) = \lVert \mathbf{c}''(s) \rVert$, the inverse of the radius of the circle that best fits
the vessel at $s$.
Two descriptors were derived from it.
The mean curvature over the first $L$\,mm of the artery is
\begin{equation}
  \bar\kappa = \frac{1}{L}\int_0^L \kappa(s)\,\mathrm{d}s,
  \label{eq:method-kappa}
\end{equation}
and the curvature profile $\kappa(s)$ on $[0, L]$ keeps where along the artery each bend lies.
Mean curvature was the geometric measure most closely related to low wall shear stress in 127
patients without coronary artery disease \cite{kashyap2022tortuosity}, and plaque concentrates in
the proximal segments of each artery \cite{araki2020plaques}, which motivates keeping the position
of a bend as well as its size.
The profile was summarized by functional principal component analysis fitted on ImageCAS-X, and the
first $K$ component scores enter the model.
% TODO: L, the rule for arteries shorter than L, and K, all fixed on ImageCAS-X.

\subsection{Features Considered and Not Used}
\label{sec:method-dropped}

The tortuosity index, the ratio of a vessel's length to the distance between its ends, is the most
reported coronary shape measure, but it showed no relation to low wall shear stress
\cite{kashyap2022tortuosity}.
Root-mean-square curvature measures the same bending as the mean with a noisier derivative.
Torsion depends on the third derivative of the centerline and could not be separated from
centerline noise on ImageCAS-X.
The Finet ratio \cite{finet2008fractal} is an exact function of $\rho$ and $\Delta_{HK}$, so its
weight in a model could not be interpreted alongside them.


\section{Reliability on the Development Dataset}
\label{sec:method-reliability}

Every feature inherits the errors of the segmentation, the centerline and the labeling, so its
sensitivity to each was measured on ImageCAS-X before the Herlev--{\O}sterbro outcomes were used.
Segmentation error was simulated by dilating and eroding the reference segmentation by one voxel.
The placement of the measurement window was varied between 0.5, 1 and 2 junction radii past the
branch point, and its length between 2, 3 and 5\,mm.
Features were also computed on trees extracted from predicted segmentations and compared with the
same features on the delivered trees.
A feature whose variation under these perturbations was large relative to its spread across
participants was dropped or replaced before modeling.
% TODO: the reliability measure (ICC?) and the threshold.

The same dataset was used to fix the smoothing of the centerline, the length $L$, the number of
principal components $K$, and the correlation structure of the features.
Fixing these choices before the outcomes are used prevents them from being tuned to the result.


\section{Statistical Analysis}
\label{sec:method-statistics}

Four models of increasing complexity were fitted for each artery domain, each adding one element to
the one before.
Model~0 is a logistic regression on age, sex, dominance and the classical risk factors, and is the
baseline the geometry has to improve on.
Model~1 adds the bifurcation features and the mean curvature $\bar\kappa$.
Model~2 replaces $\bar\kappa$ with the principal component scores of the curvature profile, and its
likelihood-ratio test against Model~1 tests whether the position of a bend adds information to
its mean.
Model~3 is a small one-dimensional convolutional network on the curvature profile combined with
the scalar features, and is exploratory, since its number of parameters exceeds what the number of
events can support.
Age, sex and dominance enter every model, as the normal value of several features depends on
them; a left-dominant tree has a larger LCx and therefore a lower normal $\lambda$.
The RCA has no bifurcation feature in this set, and the LM is too short for a curvature profile,
so these domains receive the subset of models their features allow.

The required sample size was calculated per domain with the criteria of Riley et al.
\cite{riley2020samplesize}.
% TODO: result of the calculation once the number of events is known.
Models were fitted with a ridge penalty \cite{hoerl1970ridge} chosen by nested cross-validation,
which also stabilizes the weights of correlated features.
The scan interval differs between participants and the time of plaque onset is unknown, so the
interval entered every model as a covariate.
Participants who died, had a myocardial infarction or were not rescanned were compared with the
rescanned participants on their baseline features, to assess whether the rescanned population is
selected.

Performance was estimated by bootstrap resampling of whole participants, since one participant
contributes an outcome to every domain.
The optimism-corrected area under the receiver operating characteristic curve, the calibration
slope and intercept, and the Brier score are reported, following TRIPOD+AI
\cite{collins2024tripodai}.
All validation is internal, as no second cohort with repeated CCTA in a general population was
available.
Correlations and generalized variance inflation factors \cite{fox1992gvif} between features were
reported on ImageCAS-X and on Herlev--{\O}sterbro.
The contribution of each feature family was assessed by a likelihood-ratio test, and for Model~3 by
refitting without the family rather than by permuting it \cite{hooker2021permutation}.


\section{Deviations and Implementation}
\label{sec:method-deviations}

\begin{table}[htbp]
  \centering
  \small
  \begin{tabular}{@{}llll@{}}
    \toprule
    Setting & Published & This work & Reason \\
    \midrule
    % TODO: one row per deviation from the published segmentation training settings,
    % e.g. hardware, mixed precision, resumed training across job time limits.
    \ldots & \ldots & \ldots & \ldots \\
    \bottomrule
  \end{tabular}
  \caption{Deviations from the published training settings of the segmentation model.}
  \label{tab:method-deviations}
\end{table}

% TODO: compute (cluster, GPU type, training time), software versions, and any problem found
% during the work that changed a number, written as discovery, cause, consequence, fix and check
% (style guide section 6).
